% ====================================================================
% KONFIGURASI PAKET LaTeX
% ====================================================================
% File ini berisi semua paket (packages) yang digunakan dalam dokumen.
% Paket-paket diorganisir berdasarkan fungsi untuk kemudahan pemeliharaan.
%
% JANGAN MENAMBAHKAN \documentclass DI SINI - hanya \usepackage
% ====================================================================

% --- PAKET ENCODING & BAHASA ---
% Paket untuk encoding karakter UTF-8 dan dukungan bahasa Indonesia
\usepackage[utf8]{inputenc}                 % Encoding input UTF-8
\usepackage[T1]{fontenc}                    % Encoding font
\usepackage[indonesian]{babel}              % Dukungan bahasa Indonesia

% --- PAKET FONT & FORMATASI DASAR ---
% Mengatur jenis font dan margin sesuai juknis UNEJ
\usepackage{times}                          % Font Times New Roman (sesuai juknis)
\usepackage[top=4cm, left=4cm, bottom=3cm, right=3cm]{geometry}  % Margin sesuai juknis
\usepackage{setspace}                       % Pengaturan spasi antar baris
\usepackage{listings}

% --- PAKET KUSTOMISASI VISUAL STRUKTUR ---
% Paket untuk mengatur format judul bab, section, dan elemen visual lainnya
\usepackage{titlesec}                       % Kustomisasi judul chapter/section
\usepackage{fancyhdr}                       % Header dan footer custom
\usepackage{caption}                        % Pengaturan caption gambar/tabel

% --- PAKET KONTEN & REFERENSI ---
% Paket untuk menyisipkan berbagai elemen konten
\usepackage{graphicx}                       % Menyisipkan gambar (PNG, JPG, PDF)
\usepackage{booktabs}                       % Tabel berkualitas tinggi (profesi)
\usepackage{tabularx}                       % Tabel dengan kolom yang fleksibel
\usepackage{longtable}                      % Tabel multi-halaman
\usepackage{array}                          % Format kolom tabel lanjutan
\usepackage{float}                          % Opsi penempatan float seperti [H]
\usepackage{multirow}                       % Baris multi pada tabel
\usepackage{enumitem}                       % Kustomisasi list (mis. opsi [nosep])
\usepackage{amsmath}                        % Notasi matematika advanced
\usepackage{amssymb}                        % Simbol matematika tambahan
\usepackage{natbib}                         % Manajemen referensi (APA Style)
\usepackage{lipsum}                         % Teks dummy untuk placeholder (opsional)

% --- PAKET HYPERLINK & INTERAKTIF ---
% Paket untuk membuat dokumen PDF interaktif dengan hyperlink
\usepackage{hyperref}                       % Hyperlink dalam PDF
\usepackage{tocloft}                        % Kustomisasi daftar isi/tabel/gambar
\usepackage{afterpage}                       % Menunda perintah sampai halaman berikutnya (untuk page style)
\usepackage{indentfirst}                    % Indent paragraf pertama setelah heading
\usepackage{soul}                      % Highlighting teks (opsional, untuk revisi atau catatan)

% --- PAKET TAMBAHAN (OPSIONAL) ---
% Paket tambahan yang mungkin diperlukan:
\usepackage{subcaption}
% \usepackage{listings}                     % Untuk menampilkan kode sumber
% \usepackage{xcolor}                       % Dukungan warna tambahan
% \usepackage{tikz}                         % Menggambar diagram dan grafik
% \usepackage{array}                        % Pengaturan array dan tabel lebih lanjut
% \usepackage{multicol}                     % Multiple columns layout
% \usepackage{appendix}                     % Paket untuk lampiran (sudah built-in)

% ====================================================================
% PENGATURAN SPASI GLOBAL
% Sesuai juknis: 1.5 spasi untuk isi dokumen
% ====================================================================
\setstretch{1.5}  % Atur spasi antar baris menjadi 1.5
